% -----------------------------------------------------------------------------
% This file is automatically generated.
% Do not edit manually.
% Source: notebooks/support/regression-analysis/support.Rmd
% -----------------------------------------------------------------------------


\noindent A practical way to address first-order autocorrelation is to apply an AR(1) transformation. This notebook shows the transformation explicitly rather than calling a separate named routine. We first estimate the lag-1 autocorrelation in the residuals, then quasi-difference the response and the predictors, and finally refit the model to the transformed data.
We first estimate an AR(1) model on the residuals from \ttblue{mod.3}.

\par\addvspace{\topsep}
\begingroup
\fvset{listparameters={%
  \setlength{\topsep}{0pt}%
  \setlength{\partopsep}{0pt}%
  \setlength{\parsep}{0pt}%
  \setlength{\itemsep}{0pt}%
}}
\begin{Verbatim}[breaklines=true,formatcom=\color{ada_blue}]
rho = np.corrcoef(ussteamco_mod_3.resid[1:], ussteamco_mod_3.resid[:-1])[0, 1]
print(rho)
\end{Verbatim}
\begin{Verbatim}[breaklines=true,formatcom=\color{ada_light_blue}]
      ar1
0.4413799
\end{Verbatim}
\endgroup
\par\addvspace{\topsep}
